\section{Measurable causal graphs: proofs}
\label{sec:annex-p16-measurable-causal-graphs}

\begin{annexitem}{annex:p16.measurable-ci-calculus}{thm:p16-measurable-ci-calculus}{Conditional expectation and one common measure identity}
\begin{proof}
Suppose first that the product-test definition holds. For bounded
\(u(X)\), put \(h(Z)=E[u(X)\mid Z]\). On a generating event
\(\{Y\in D,Z\in E\}\), conditional pullout and the definition give
\[
\begin{aligned}
 E[\mathbf1_E(Z)\mathbf1_D(Y)h(Z)]
 &=E[\mathbf1_E(Z)h(Z)E[\mathbf1_D(Y)\mid Z]]\\
 &=E[\mathbf1_E(Z)u(X)\mathbf1_D(Y)].
\end{aligned}
\]
A pi-lambda argument extends this integral identity to all events in
\(\sigma(Y,Z)\), proving that \(h(Z)\) is a conditional expectation
given \((Y,Z)\). Conversely if the conditional expectation given
\((Y,Z)\) is \(h(Z)\), tower and pullout give
\[
 E[u(X)v(Y)\mid Z]
 =E[h(Z)v(Y)\mid Z]=h(Z)E[v(Y)\mid Z].
\]
The requisite tower and pullout results are
\cref{thm:p15-conditional-tower,thm:p15-conditional-pullout}.

Choose regular kernels for \(X,Y,(X,Y)\) given \(Z\). Their existence
and uniqueness as measures follow from
\cref{thm:p15-value-conditional-law}. Each standard Borel target has
a countable generating pi-system, by its Borel real coding. For every
rectangle from the two countable systems, conditional independence of
its indicators equates the joint kernel to the product of its marginal
kernels outside a null set of \(Z\). Remove the countable union of these
sets. At every remaining \(z\), measure uniqueness gives equality on
the entire product sigma-algebra. The product kernel is measurable:
its values on rectangles are products of measurable marginal values,
and a pi-lambda argument extends this to all measurable sets. Conversely
the common measure identity integrated against products of bounded
functions gives conditional independence. This proves the second
equivalence without separate null-set choices for uncountably many events.

Symmetry is immediate from the defining product identity. Decomposition
follows by taking a test function that ignores an extra coordinate.
For weak union, the no-extra-information identity gives
\(E[u(X)\mid Y,W,Z]=E[u(X)\mid Z]\). The right side is measurable
from \((W,Z)\); tower makes it also \(E[u(X)\mid W,Z]\).
Applying the equivalence in reverse gives \(X\perp Y\mid(W,Z)\).
For contraction the two premises give successively
\[
 E[u(X)\mid Y,W,Z]=E[u(X)\mid Y,Z]=E[u(X)\mid Z],
\]
which is precisely the conclusion. All tests are bounded, so every
conditional expectation used is integrable. Measurable transformations
of independent tuples follow either by composing bounded tests or,
when the functions also depend on \(Z\), by pushing forward the
conditional product kernels at each \(z\). No intersection axiom or
strictly positive law entered any of these steps.
\end{proof}
\end{annexitem}

\begin{annexitem}{annex:p16.measurable-kernel-construction}{thm:p16-measurable-kernel-construction}{Associative kernel combination and legal adjacent swaps}
\begin{proof}
For the empty graph its point-mass law gives all assertions directly.
At the first vertex of a nonempty graph the parent tuple is empty, so its kernel is a
probability measure. At each later stage extend its kernel to the whole
preceding tuple by the measurable parent projection. The kernel
joint-measure theorem constructs a unique new probability measure and
identifies its conditional law of the new coordinate. Induction gives
existence, normalization, the iterated-integral identity, and the
conditional-prefix statement. All parameter integrals are measurable by
the same theorem. Nonnegative tests pass through each step by monotone
convergence; signed tests are handled only after their absolute integrals
are finite. In particular no multiplication of hypothetical densities
has occurred.

For order independence, take two consecutive vertices \(v,w\) in a
topological order with neither a parent nor an ancestor of the other.
Both parent tuples are functions of the common earlier tuple \(z\).
Conditional on \(z\), their two update kernels combine as the product
\(K_v(z,dx_v)\otimes K_w(z,dx_w)\). On rectangles this follows
directly from the kernel joint-measure definition; measure uniqueness
extends it to every measurable set. Fubini for these two probability
measures permits either integration order. Integrating all later
updates first gives a measurable nonnegative function of \((z,x_v,x_w)\),
so the same swap preserves the full law, not just this pair's marginal.

Any two topological orders are connected by such adjacent swaps. To see
this, take the first vertex in the desired order and move it leftward
through the current order. A crossed vertex cannot be its ancestor,
since it would then have to precede it in the desired order; it cannot
be its descendant, since it precedes it in the current order. Thus each
swap is of incomparable vertices and is legal. Fix that first position
and repeat on the remaining finite graph. Induction proves the claim
and hence order independence. This is the source's kernel-combination
argument with the needed product/Fubini domains made explicit.

If \(D\) is ancestral, all parents of its vertices lie in \(D\).
List a topological order of \(G[D]\) first, followed by one of the
remaining vertices. No edge can enter \(D\) from its complement,
so this is a topological order of \(G\). Its initial joint law is
the kernel construction on \(D\), and all later extensions preserve
that marginal. Order independence now proves the asserted ancestral
marginal for every original ordering. Topological prefixes are
ancestral. Their conditional-coordinate identities therefore hold
also in the full law, since conditional expectations from preceding
coordinates depend only on the corresponding joint marginal.
\end{proof}
\end{annexitem}

\begin{annexitem}{annex:p16.measurable-kernel-randomization}{thm:p16-measurable-kernel-randomization}{A jointly measurable quantile and Borel image repair}
\begin{proof}
Choose a measurable isomorphism \(\phi:S\to E\), where \(E\) is a
nonempty Borel subset of \(\mathbb R\), using the proved standard
Borel coding. Push each \(K(t,\cdot)\) through \(\phi\) and regard
it as a probability measure \(\widetilde K_t\) on \(\mathbb R\)
concentrated on \(E\). Its CDF is
\(F_t(x)=\widetilde K_t(( -\infty,x])\). For fixed \(x\), this
is measurable in \(t\). Define the source's upper quantile by the
equivalent rational formula
\[
 q(t,u)=\inf\{r\in\mathbb Q:F_t(r)>u\},\qquad 0<u<1.
\]
For each \((t,u)\), the set is nonempty by the upper CDF limit and
is bounded below by the lower CDF limit, so the infimum is real.
Its strict sublevel sets satisfy
\[
 \{(t,u):q(t,u)<x\}
   =\bigcup_{r\in\mathbb Q,\ r<x}\{(t,u):F_t(r)>u\}.
\]
This countable union is measurable, proving joint measurability.

For fixed \(t,x\), monotonicity and right continuity of the CDF give
\[
 \{u:u<F_t(x)\}\subseteq\{u:q(t,u)\leq x\}
                  \subseteq\{u:u\leq F_t(x)\}.
\]
For the first inclusion use rational arguments arbitrarily close above
\(x\). For the second, if the infimum is at most \(x\), every rational
argument strictly above \(x\) has CDF at least \(u\), and then let
such arguments decrease to \(x\). The two bounding sets have the same
uniform measure \(F_t(x)\). Thus \(q(t,U)\) has CDF \(F_t\), and
has law \(\widetilde K_t\). In particular \(q(t,U)\in E\) almost
surely for each fixed \(t\).

Choose \(s_0\in S\) and define
\[
 F(t,u)=\begin{cases}\phi^{-1}(q(t,u)),&q(t,u)\in E,\\
                         s_0,&q(t,u)\notin E.\end{cases}
\]
The exceptional set is jointly measurable because \(E\) is Borel.
The inverse coding and this measurable default make \(F\) jointly
measurable on its whole domain. For each fixed parameter the modification
is on a uniform-null set, and so its pushforward law is exactly
\(K(t,\cdot)\). This proves the kernel realization at every parameter,
including parameters absent from an observed law.

Apply this construction separately to each of the finitely many DAG
kernels and use independent uniform variables on their finite product
space. Topological evaluation and successive conditioning show that its
joint law is the kernel law. The representation is not unique: a fair
response independent of binary treatment can be represented as
\(Y=U\), or as \(Y=A\oplus U\), with a fair independent bit \(U\).
Both give the same controlled response kernel at both actions, but
their potential responses are respectively equal or opposite. Thus
the construction provides a structural realization, not an inference
of a unique cross-world coupling from a probability kernel.
\end{proof}
\end{annexitem}

\begin{annexitem}{annex:p16.measurable-global-markov}{thm:p16-measurable-global-markov}{The complete density-free terminal-vertex induction}
\begin{proof}
\paragraph{Recursive kernels imply the global property.}
We first prove a full-partition assertion by induction on the number
of vertices. Suppose \(N=D\mathbin{\dot\cup}E\mathbin{\dot\cup}C\)
and there is no edge between \(D\) and \(E\) in the moral graph
after deleting \(C\). We claim \(X_D\perp X_E\mid X_C\).
It is immediate if either side is empty, including the empty and
one-vertex starting cases.

Take the last vertex \(v\) in a topological order. It is terminal.
The remaining law is its kernel-model marginal, and the moral graph
on the remaining DAG has no more edges than the old moral graph
with \(v\) deleted. If \(v\in D\), its parents lie in
\((D\setminus\{v\})\cup C\), since a parent in \(E\) would be
a forbidden moral edge. The conditional kernel of \(X_v\) given
all other variables consequently depends only on
\(X_{D\setminus\{v\}},X_C\). The no-extra-information
characterization gives
\(X_v\perp X_E\mid X_{D\setminus\{v\}},X_C\).
Induction gives \(X_{D\setminus\{v\}}\perp X_E\mid X_C\).
Symmetry and contraction combine these into the desired independence
of \(X_D\) and \(X_E\). The case \(v\in E\) is symmetric.

If \(v\in C\), put \(C'=C\setminus\{v\}\). Parents of a
common child form a clique in the moral graph. Thus its parents
outside \(C'\) cannot meet both \(D\) and \(E\). Suppose they
are all in \(D\); the other case is symmetric. The terminal kernel
gives \(X_v\perp X_E\mid X_D,X_{C'}\), while induction on the
remaining graph gives \(X_D\perp X_E\mid X_{C'}\). Symmetry and
contraction give \((X_D,X_v)\perp X_E\mid X_{C'}\); weak union
then gives \(X_D\perp X_E\mid X_{C'},X_v\). This is the required
conditioning on \(X_C\). The argument also covers the case where
all parents lie in \(C'\). This completes the full-partition induction.

Now let disjoint \(A,B,C\) be d-separated in \(G\). Restrict to
\(H=\operatorname{An}_G(A\cup B\cup C)\). The ancestral marginal
lemma retains its kernels. By the graph-theoretic moral bridge, no
component of the moral graph of \(G[H]\) after deletion of \(C\)
meets both \(A\) and \(B\). Assign all components meeting \(A\)
to \(D\), and all remaining components to \(E\). This fills in
the vertices not originally named in the query and gives a full
partition of \(H\) with no \(D\)-to-\(E\) edge, with
\(A\subseteq D,B\subseteq E\). The induction just proved gives
\(X_D\perp X_E\mid X_C\); decomposition gives the desired query.
Empty query sets are immediate. No density, positivity or intersection
property was used.

\paragraph{Global implies local, and local implies ordered.}
For a vertex \(v\), the set consisting of \(v\) and its
nondescendants is ancestral. Within it \(v\) has no children.
Its only moral neighbors are its parents: any other moral neighbor
would require a common child, which would be a descendant of \(v\).
Deleting the parents therefore separates \(v\) from its other
nondescendants. The global property gives the local property.
Every predecessor is a nondescendant, so decomposition gives the
ordered property.

\paragraph{Ordered implies recursive kernels.}
For each vertex choose a regular conditional kernel of \(X_v\)
given its parents, specified by a measurable probability default on
any null parent arguments. The ordered property and the first
conditional-independence equivalence make this same kernel a
conditional law of \(X_v\) given all predecessors. In detail,
compare its values on a countable generating class with a regular
law given all predecessors; outside their countable null union they
agree as measures. The joint-measure identity for conditional kernels
then expresses each successive actual prefix law as its previous
prefix combined with the chosen kernel. Induction gives the full
recursive kernel law. Together the implications prove all four
equivalences, with the exact standard Borel existence hypotheses.
\end{proof}
\end{annexitem}

\begin{annexitem}{annex:p15-w02b-p15-009-sg-trivial}{thm:p15-w02b-p15-009-semigraphoid-trivial}{Empty sides by contraction}
\begin{proof}
Throughout, an application of an axiom names the sets placed in its
slots \((A,B,C,D)\); the slots are mutually disjoint in every
application, and any of them may be empty.

List \(V\) in the given topological order as \(v_1<\dots<v_n\), where
\(n\geq1\), and put \(P_k=\{v_1,\ldots,v_k\}\), so that
\(\operatorname{pr}(v_k)=P_{k-1}\) and \(P_0=\varnothing\). Fix \(k\)
and write \(R_k=\operatorname{pr}(v_k)\setminus\operatorname{pa}(v_k)\).
Every parent precedes its child, so
\(\operatorname{pa}(v_k)\subseteq P_{k-1}\). Property (O) gives
\(v_k\perp_\sigma R_k\mid\operatorname{pa}(v_k)\). Weak union (S3)
with slots \((\{v_k\},\varnothing,\operatorname{pa}(v_k),R_k)\) turns
this into \(v_k\perp_\sigma\varnothing\mid\operatorname{pa}(v_k)\cup R_k\),
that is \(v_k\perp_\sigma\varnothing\mid P_{k-1}\). Symmetry (S1)
gives
\[
 \varnothing\perp_\sigma\{v_k\}\mid P_{k-1}\qquad(1\leq k\leq n).
\]

We show \(\varnothing\perp_\sigma P_k\mid\varnothing\) by induction on
\(k\). For \(k=1\) it is the display with \(P_0=\varnothing\). If it
holds for \(k-1\), contraction (S4) with slots
\((\varnothing,P_{k-1},\varnothing,\{v_k\})\) has premises
\(\varnothing\perp_\sigma P_{k-1}\mid\varnothing\) and
\(\varnothing\perp_\sigma\{v_k\}\mid P_{k-1}\), and its conclusion is
\(\varnothing\perp_\sigma P_k\mid\varnothing\). At \(k=n\) this is
\(\varnothing\perp_\sigma V\mid\varnothing\).

Now let \(A,C\subseteq V\) be disjoint and put
\(R=V\setminus(A\cup C)\). Decomposition (S2) with slots
\((\varnothing,A\cup C,\varnothing,R)\) gives
\(\varnothing\perp_\sigma A\cup C\mid\varnothing\). Weak union (S3)
with slots \((\varnothing,A,\varnothing,C)\) gives
\(\varnothing\perp_\sigma A\mid C\), and symmetry (S1) gives
\(A\perp_\sigma\varnothing\mid C\). The hypothesis \(n\geq1\) was used
to start the induction; the warning in the body shows it is needed.
\end{proof}
\end{annexitem}

\begin{annexitem}{annex:p15-w02b-p15-009-sg-markov}{thm:p15-w02b-p15-009-semigraphoid-markov}{Induction over ancestral sets}
\begin{proof}
Axiom applications name their slots \((A,B,C,D)\) as in the previous
proof. For \(H\subseteq V\), let \(\mathsf m(H)\) be the undirected
graph on \(H\) in which distinct \(i,j\) are adjacent when they are
joined by an edge of \(D\) or are both parents of some \(k\in H\).
When \(H=\operatorname{an}_D(A\cup B\cup S)\) this is the moral graph
\(D_H^{\mathrm m}\) of \cref{def:p16-ancestral-moral-graph}.

\paragraph{Four graph facts.}
(a) \emph{For every \(\alpha\), the set \(\{\alpha\}\cup\operatorname{nd}(\alpha)\)
is ancestral.} Let \(u\) have a directed path of positive length to
some \(w\in\{\alpha\}\cup\operatorname{nd}(\alpha)\). If
\(u\in\operatorname{de}(\alpha)\), joining a path from \(\alpha\) to
\(u\) with the path from \(u\) to \(w\) gives \(w\in\operatorname{de}(\alpha)\).
That is impossible: \(\operatorname{nd}(\alpha)\) misses
\(\operatorname{de}(\alpha)\) by definition, and
\(\alpha\notin\operatorname{de}(\alpha)\) by acyclicity. Hence
\(u\in\{\alpha\}\cup\operatorname{nd}(\alpha)\).

(b) \emph{\(\operatorname{pa}(\alpha)\subseteq\operatorname{pr}(\alpha)
\subseteq\operatorname{nd}(\alpha)\).} Parents precede their child.
A predecessor \(\beta\) differs from \(\alpha\), and a directed path
from \(\alpha\) to \(\beta\) would force \(\alpha<\beta\).

(c) \emph{If \(H\) is ancestral and \(v\in H\) has no child in \(H\),
then \(H\setminus\{v\}\subseteq\operatorname{nd}(v)\) and
\(H\setminus\{v\}\) is ancestral.} If \(w\in H\setminus\{v\}\) were in
\(\operatorname{de}(v)\), the first vertex after \(v\) on a directed
path from \(v\) to \(w\) would be an ancestor of \(w\) or \(w\) itself,
hence in \(H\), and a child of \(v\), contrary to the hypothesis. For the second claim, an ancestor
of a vertex of \(H\setminus\{v\}\) lies in \(H\); it is not \(v\),
since the first step of its path from \(v\) would be a child of \(v\)
in \(H\).

(d) \emph{If \(H'\subseteq H\), every edge of \(\mathsf m(H')\) is an
edge of \(\mathsf m(H)\).} An edge of \(D\) between vertices of \(H'\)
joins vertices of \(H\), and a common child in \(H'\) lies in \(H\).

\paragraph{(G) implies (L).}
Fix \(\alpha\) and put \(A=\{\alpha\}\),
\(B=\operatorname{nd}(\alpha)\setminus\operatorname{pa}(\alpha)\) and
\(S=\operatorname{pa}(\alpha)\). They are pairwise disjoint by (b).
By (a) and since \(A\cup B\cup S=\{\alpha\}\cup\operatorname{nd}(\alpha)\),
the set \(H=\operatorname{an}_D(A\cup B\cup S)\) equals
\(\{\alpha\}\cup\operatorname{nd}(\alpha)\). A neighbour \(w\) of
\(\alpha\) in \(\mathsf m(H)\) is a parent of \(\alpha\): a child of
\(\alpha\) is a descendant, so not in \(H\), and a common child of
\(\alpha\) and \(w\) would be a child of \(\alpha\) in \(H\). After
deleting \(S\), the vertex \(\alpha\) is therefore isolated, and no path
joins \(A\) to \(B\). By \cref{thm:p16-moral-separation-bridge},
\(A\mathrel{\perp_D}B\mid S\), and (G) gives the local statement at
\(\alpha\). This uses no hypothesis on \(V\).

\paragraph{(L) implies (O).}
By (b), \(\operatorname{nd}(\alpha)\setminus\operatorname{pa}(\alpha)\)
is the disjoint union of
\(\operatorname{pr}(\alpha)\setminus\operatorname{pa}(\alpha)\) and
\(R=\operatorname{nd}(\alpha)\setminus\operatorname{pr}(\alpha)\).
Decomposition (S2) with slots
\((\{\alpha\},\operatorname{pr}(\alpha)\setminus\operatorname{pa}(\alpha),
\operatorname{pa}(\alpha),R)\) gives (O) at \(\alpha\), for every
topological order.

\paragraph{(O) implies (L).}
We prove, by induction on \(n=|V|\), that for every DAG on \(n\)
vertices, every topological order and every semigraphoid, (O) implies
(L). For \(n=0\) both are vacuous. Let \(n\geq1\), let \(v^*\) be the
last vertex in the order, and put \(V'=V\setminus\{v^*\}\) and
\(D'=D[V']\). The restriction of \(\perp_\sigma\) to triples of subsets
of \(V'\) is a semigraphoid, since every axiom's conclusion uses only
unions of its premises' sets. Every child follows its parent, so
\(v^*\) has no child; therefore for \(\alpha\in V'\) its parents,
predecessors and descendants other than \(v^*\) are the same in \(D'\)
as in \(D\), and \(\operatorname{nd}_{D'}(\alpha)=
\operatorname{nd}(\alpha)\setminus\{v^*\}\). Thus (O) for \(D'\) and
the restricted order consists of statements of (O) for \(D\), and the
induction hypothesis gives
\[
 \alpha\perp_\sigma\bigl(\operatorname{nd}(\alpha)\setminus\{v^*\}\bigr)
   \setminus\operatorname{pa}(\alpha)\mid\operatorname{pa}(\alpha)
 \qquad(\alpha\in V').
\]
Fix \(\alpha\in V'\). If \(v^*\in\operatorname{de}(\alpha)\), this is
(L) at \(\alpha\). Otherwise \(v^*\in\operatorname{nd}(\alpha)\). Put
\(N=\operatorname{nd}(\alpha)\setminus\{v^*\}\). No parent of \(v^*\)
is \(\alpha\) or a descendant of \(\alpha\), since then \(v^*\) would be
a descendant; so \(\operatorname{pa}(v^*)\subseteq N\). Since
\(\operatorname{pr}(v^*)=V'\), property (O) at \(v^*\) reads
\(v^*\perp_\sigma V'\setminus\operatorname{pa}(v^*)\mid\operatorname{pa}(v^*)\).
The set \(\{\alpha\}\cup(N\setminus\operatorname{pa}(v^*))\) is contained
in \(V'\setminus\operatorname{pa}(v^*)\), so decomposition (S2) gives
\(v^*\perp_\sigma\{\alpha\}\cup(N\setminus\operatorname{pa}(v^*))\mid
\operatorname{pa}(v^*)\). Weak union (S3) with slots
\((\{v^*\},\{\alpha\},\operatorname{pa}(v^*),N\setminus\operatorname{pa}(v^*))\)
gives \(v^*\perp_\sigma\alpha\mid N\), and symmetry (S1) gives
\(\alpha\perp_\sigma v^*\mid N\). Now
\(\operatorname{pa}(\alpha)\subseteq N\) by (b), and \(v^*\) is not a
parent of \(\alpha\). Contraction (S4) with slots
\((\{\alpha\},N\setminus\operatorname{pa}(\alpha),\operatorname{pa}(\alpha),\{v^*\})\)
combines the displayed statement with \(\alpha\perp_\sigma v^*\mid N\)
into \(\alpha\perp_\sigma\operatorname{nd}(\alpha)\setminus\operatorname{pa}(\alpha)
\mid\operatorname{pa}(\alpha)\). Finally, at \(\alpha=v^*\) the vertex
has no descendants, so
\(\operatorname{nd}(v^*)=V'=\operatorname{pr}(v^*)\) and (L) at \(v^*\)
is (O) at \(v^*\). This completes the induction. It uses no hypothesis
on \(V\).

\paragraph{(L) implies (G) when \(V\neq\varnothing\).}
Assume (L). By the previous paragraphs (O) holds, so
\cref{thm:p15-w02b-p15-009-semigraphoid-trivial} gives every statement
with an empty side. We prove by induction on \(|H|\) the following
assertion for every ancestral \(H\subseteq V\):
\begin{quote}
if \(H\) is the disjoint union of \(P\), \(Q\) and \(C\), and no edge
of \(\mathsf m(H)\) joins \(P\) to \(Q\), then \(P\perp_\sigma Q\mid C\).
\end{quote}
If \(P\) or \(Q\) is empty, the statement has an empty side and holds.
This covers \(H=\varnothing\). Let \(P,Q\) be nonempty, so
\(H\neq\varnothing\), and let \(v\) be the last vertex of \(H\) in the
order. A child of \(v\) follows \(v\), so \(v\) has no child in \(H\).
By (c), \(H'=H\setminus\{v\}\) is ancestral and
\(H'\subseteq\operatorname{nd}(v)\). By (d), \(\mathsf m(H')\) has no
edge between \(P\setminus\{v\}\) and \(Q\setminus\{v\}\). The parents
of \(v\) lie in \(H'\).

\emph{Case \(v\in X\), where \(\{X,Y\}=\{P,Q\}\).} A parent of \(v\) in
\(Y\) would give an edge of \(\mathsf m(H)\) between \(P\) and \(Q\),
so \(\operatorname{pa}(v)\subseteq(X\setminus\{v\})\cup C\). Put
\(R=((X\setminus\{v\})\cup C)\setminus\operatorname{pa}(v)\). Then
\(Y\cup R\subseteq H'\setminus\operatorname{pa}(v)
\subseteq\operatorname{nd}(v)\setminus\operatorname{pa}(v)\).
From (L) at \(v\), decomposition (S2) gives
\(v\perp_\sigma Y\cup R\mid\operatorname{pa}(v)\), and weak union (S3)
with slots \((\{v\},Y,\operatorname{pa}(v),R)\) gives
\(v\perp_\sigma Y\mid(X\setminus\{v\})\cup C\). The induction
hypothesis for \(H'\), partitioned into \(X\setminus\{v\}\), \(Y\) and
\(C\), gives \(X\setminus\{v\}\perp_\sigma Y\mid C\). After symmetry
(S1) on both statements, contraction (S4) with slots
\((Y,X\setminus\{v\},C,\{v\})\) gives \(Y\perp_\sigma X\mid C\).
Symmetry (S1) then gives \(X\perp_\sigma Y\mid C\), and with
\(\{X,Y\}=\{P,Q\}\) one of these two is \(P\perp_\sigma Q\mid C\).

\emph{Case \(v\in C\).} Put \(C'=C\setminus\{v\}\). The parents of
\(v\) outside \(C'\) lie in \(P\cup Q\) and are pairwise adjacent in
\(\mathsf m(H)\), as parents of the common child \(v\in H\). So they
cannot meet both \(P\) and \(Q\): choose \(\{X,Y\}=\{P,Q\}\) with
\(\operatorname{pa}(v)\subseteq X\cup C'\). Put
\(R=(X\cup C')\setminus\operatorname{pa}(v)\). Then
\(Y\cup R\subseteq\operatorname{nd}(v)\setminus\operatorname{pa}(v)\),
and as in the first case (L), (S2) and (S3) give
\(v\perp_\sigma Y\mid X\cup C'\); symmetry gives
\(Y\perp_\sigma v\mid X\cup C'\). The induction hypothesis for \(H'\),
partitioned into \(X\), \(Y\) and \(C'\), gives
\(X\perp_\sigma Y\mid C'\), hence \(Y\perp_\sigma X\mid C'\) by (S1).
Contraction (S4) with slots \((Y,X,C',\{v\})\) gives
\(Y\perp_\sigma X\cup\{v\}\mid C'\). Weak union (S3) with slots
\((Y,X,C',\{v\})\) gives \(Y\perp_\sigma X\mid C\), and symmetry gives
\(X\perp_\sigma Y\mid C\). Again one of the two is
\(P\perp_\sigma Q\mid C\). This completes the induction.

Now let \(A,B,S\) be pairwise disjoint with \(A\mathrel{\perp_D}B\mid S\).
If \(A\) or \(B\) is empty, the lemma gives \(A\perp_\sigma B\mid S\).
Otherwise put \(H=\operatorname{an}_D(A\cup B\cup S)\); it is ancestral,
since an ancestor of an ancestor is an ancestor. By
\cref{thm:p16-moral-separation-bridge}, deleting \(S\) from
\(\mathsf m(H)\) separates \(A\) from \(B\). Let \(P\) be the union of
the vertex sets of those components of \(\mathsf m(H)\setminus S\) that
meet \(A\), and put \(Q=H\setminus(P\cup S)\). Then \(A\subseteq P\),
and \(B\subseteq Q\) because no component meets both \(A\) and \(B\).
An edge of \(\mathsf m(H)\) from \(u\in P\) to \(w\in Q\) avoids \(S\),
so it lies in \(\mathsf m(H)\setminus S\) and puts \(w\) in the
component of \(u\), that is in \(P\); hence there is none. The
assertion gives \(P\perp_\sigma Q\mid S\). Decomposition (S2) with
slots \((P,B,S,Q\setminus B)\) gives \(P\perp_\sigma B\mid S\); symmetry
gives \(B\perp_\sigma P\mid S\); decomposition with slots
\((B,A,S,P\setminus A)\) gives \(B\perp_\sigma A\mid S\); and symmetry
gives \(A\perp_\sigma B\mid S\). This is (G).

\paragraph{Relation to the source's argument.}
The source proves (G)\(\Rightarrow\)(L) by the same isolation of
\(\alpha\), (L)\(\Rightarrow\)(O) by decomposition, (O)\(\Rightarrow\)(L)
by removing the last vertex, and (L)\(\Rightarrow\)(G) by removing a
terminal vertex \cite[Theorem~2.44 and its proof, pp.~65--66]{Lauritzen2019GraphicalLectures2020}.
In the last direction the source works with the query sets themselves.
The proof above first enlarges the query to a partition of the ancestral
set. That step is what guarantees that the parents of a conditioned
terminal vertex lie on one side, since a parent outside the three query
sets would otherwise be unassigned. It also writes out the closing
contraction in the case \(v\in X\), and treats the statements with an
empty side through the lemma, which is where \(V\neq\varnothing\) enters.
\end{proof}
\end{annexitem}

\begin{annexitem}{annex:p15-w02b-p15-009-sg-order}{thm:p15-w02b-p15-009-order-invariance}{Through the local property}
\begin{proof}
Let (O) hold relative to a topological order \(<\). The paragraph
``(O) implies (L)'' in the proof of
\cref{thm:p15-w02b-p15-009-semigraphoid-markov} gives (L), with no
hypothesis on \(V\). Property (L) does not mention an order. For any
other topological order \(<'\), the paragraph ``(L) implies (O)'' of
the same proof, applied with \(<'\), gives (O) relative to \(<'\).
This is the source's argument, routed through (L) alone, so the
nonemptiness needed for (G) is not used
\cite[Corollary~2.45, p.~66]{Lauritzen2019GraphicalLectures2020}.
\end{proof}
\end{annexitem}

\begin{annexitem}{annex:p16.measurable-interventions}{thm:p16-measurable-interventions}{Independent disturbances and measurable equation replacement}
\begin{proof}
Topological induction gives a unique solution at every background
value: each total function is evaluated after its parent values are
known. The same induction, using product measurability of finite
tuples and composition, proves joint measurability in all fixed
intervention values and background coordinates. This is the argument
of \cref{thm:p16-mediation-measurability} with the intervention
parameters retained explicitly.

For an unmodified vertex define
\(K_v(x_{\operatorname{pa}(v)},B)=
\int\mathbf1_B(f_v(x_{\operatorname{pa}(v)},u))\lambda_v(du)\),
where \(\lambda_v\) is its disturbance law. The kernel section
theorem makes this measurable in the parent tuple, and its values
are probability measures. Every predecessor in a topological order
is a function of earlier disturbances, independent of the current
\(U_v\). Successive integration therefore gives exactly this
conditional kernel. Under a constant intervention its output is
the prescribed value and its kernel is \(\delta_{t_v}\). Thus
the full intervened law is the claimed recursive kernel law.
Integrating a measurable indicator of the joint solution against
the fixed product disturbance law makes \(t\mapsto P^t(B)\)
measurable for every target event, proving a probability kernel in
the intervention parameter.

A history-based replacement kernel is realized with a new independent
uniform variable by the randomization theorem. Its parents are the
specified earlier coordinates; with the assumed acyclic augmented
graph, the same argument gives its kernel law and global Markov
property. For consistency, the factual solution already satisfies
all replacement equations when the replaced values are its attained
values, and uniqueness preserves the solution. The same reasoning
in an already modified model proves composition. These are pointwise
equation identities and require no observational support condition.
Changing a kernel on an unobserved parent argument, however, changes
the specified causal model and can change interventions there.
\end{proof}
\end{annexitem}

\begin{annexitem}{annex:p16.measurable-response-component}{thm:p16-measurable-response-component}{A whole response background on one side of the separator}
\begin{proof}
Add the private disturbance roots and take the ancestral graph of
\(Y,B,C\), including those roots, then moralize and delete \(C\).
Let \(D\) be the component containing \(Y\). It contains no vertex
of \(B\), by the separation hypothesis. Every parent of a vertex
in the ancestral graph is also in that graph. If a nonconditioning
parent has a child in \(D\), its direct edge remains in the moral
graph and hence it too lies in \(D\). In particular every private
disturbance needed by an endogenous vertex of \(D\) is among the
roots in \(R\). Parents outside \(D\) can only be in \(C\).

Set these external conditioning inputs to an arbitrary formal tuple
\(c\), and the private roots of \(D\) to an arbitrary \(r\).
Evaluate the vertices of \(D\) in their inherited topological order,
using the given parameter \(t\). Total functions define every
coordinate even if \((c,r)\) is not naturally attainable. Joint
measurability follows by induction. Denote the resulting response
function by \(F(t,c,r)\). At an actual outcome, the conditioning
tuple is unchanged by hypothesis, and all parent and disturbance
inputs agree with those in this restricted evaluation. Induction
therefore proves \(Y^t=F(t,C,R)\) simultaneously, pointwise for
every parameter. No almost-everywhere choice was involved.

Choose one parameter value \(t_0\). In the model at \(t_0\),
the tuple of roots \(R\) is d-separated from \(B\) given \(C\):
their relevant ancestral moral graph is a subgraph of the one already
considered, and its root vertices lie on the \(D\) side. The global
Markov theorem therefore gives \(R\perp B\mid C\). Both these
tuples and their background law are the same for all parameters,
so this is one assertion about their original joint law, not a
parameter-dependent collection of versions. The common product-kernel
characterization supplies one conditional measure identity.

If the random parameter is a measurable function of \((B,C)\), its
conditional joint law with \(R\) is the corresponding pushforward
of this product kernel. Integrating the jointly measurable \(F\)
against that joint law proves the stated random-substitution license,
for bounded tests or with explicit absolute integrability for means.
This is exactly the uniform premise unavailable from separate
pointwise conditional-response assertions.
\end{proof}
\end{annexitem}

\begin{annexitem}{annex:p16.measurable-mediation-fields}{thm:p16-measurable-mediation-fields}{Three uniform response fields from three different graph separations}
\begin{proof}
All the graphs in this proof are finite, and all their functions remain
jointly measurable. Expose independent disturbances as private roots;
this does not change any separation among the original vertices, since
a new root has just one child. Baseline nondescendant conditions ensure
that the tuples named below are unchanged pointwise by the specified
input replacements.

\paragraph{Cross-world field.}
Take \(H=G_{\underline A\underline M}\), restrict to the ancestors
of \(\{Y\}\cup M\cup C\), moralize and remove \(C\). Let \(D_Y\)
be the component of \(Y\) and \(D_M\) the union of components meeting
the mediator set. The first source separation makes them disjoint.
Restore every original outgoing mediator edge whose child is an ancestor
of some mediator in the graph with exposure inputs cut. This restores
the mediator-to-mediator ancestral equations needed for the entire tuple
\(M_b\), while keeping the separation intact. Here are the exact
graph facts used in that claim. From the child of any restored edge,
follow a directed route to its first mediator endpoint. The route already
exists in \(H\), as it uses no outgoing mediator edge before that
endpoint. It avoids \(C\), which has no mediator descendants. Thus
the child lies in \(D_M\). Every other remaining nonbaseline parent
of that child either has an edge already in \(H\), or is itself a
mediator; it too lies in \(D_M\). New moral edges therefore join only
\(D_M\) and possibly deleted baseline vertices. All original ancestors
of the mediator set already belong to the ancestral graph, by stopping
their directed routes at the first mediator. No new ancestral vertex
is required and no new moral connection to \(D_Y\) appears.

Fix \(b\). In the restored graph, replace every exposure input to a
child in \(D_Y\) by \(a\) and every other exposure input by \(b\).
Replace every mediator input not restored by its component of \(m\).
The mediator tuple in this auxiliary model is exactly \(M_b\), by
topological evaluation on its restored ancestors. The conditioning
tuple is the original \(C\). Both are unchanged as \((a,m)\) varies.
On ancestral routes to \(Y\), stop at \(C\). All other vertices
needed there belong to \(D_Y\), and their exposure and mediator inputs
are precisely \(a,m\); hence the response is \(Y_{a,m}\) pointwise.
These equation comparisons use no value-space cardinality assumption.

Apply the response-component lemma to this parameterized auxiliary
model, with parameter \((a,m)\), unchanged tuple \(B=M_b\) and
conditioning \(C\). It produces the actual private-root tuple \(R_0\),
a jointly measurable \(F_0\), and \(R_0\perp M_b\mid C\). The
response component and its equations depend on \(a,m,C,R_0\), not
on \(b\). Thus the same \(R_0,F_0\) works for each reference value
\(b\). This establishes a common product-kernel premise for the whole
response family, rather than only independence of its separately
evaluated fixed-mediator coordinates.

\paragraph{Exposure-adjustment field.}
In \(G_{\underline A\overline M}\), retain the natural exposure
equation, feed \(a\) to its outgoing inputs, and fix all mediator
equations at \(m\). Its response is \(Y_{a,m}\). Its tuples
\(A\) and \(S=(C,W_1)\) are unchanged pointwise: every mediator
is a descendant of \(A\), and \(S\) has no exposure descendants.
The second separation and the response-component lemma, now with
parameter \((a,m)\), unchanged tuple \(A\) and conditioning \(S\),
give \(Y_{a,m}=F_1(a,m,S,R_1)\) for all parameters and
\(R_1\perp A\mid S\). Disturbances at fixed mediator coordinates
can be retained as unused inputs; their extra private roots do not
create a path between old vertices and can be ignored by \(F_1\).

\paragraph{Mediator-adjustment field.}
In this field we condition on the tuple
\(T=(A,C,W_1,W_2)\), and the graph
\(G\) carries its exposed private disturbance roots. The ancestor
operator \(\operatorname{an}\) includes its starting set, as in
\cref{def:p16-active-path}. Three facts about \(T\) are used
repeatedly. First, no coordinate of \(T\) is a descendant of a
mediator. The set \(W_2\) contains none by hypothesis. Every mediator
is a descendant of \(A\), so every mediator descendant is a descendant
of \(A\), and \(C=W_0\) and \(W_1\) contain no descendant of \(A\).
Finally \(A\) is not a mediator descendant, because every mediator
descends from \(A\) and \(G\) is acyclic. Second, \(T\) meets neither
\(M\) nor \(\{Y\}\): the covariate sets exclude \(A\), \(Y\) and the
mediators, and \(A\) is distinct from \(Y\) and from every mediator.
Third, \(Y\) is not an ancestor of any mediator, since each mediator
is an ancestor of \(Y\) and \(G\) is acyclic.

\emph{The restored graph.} Put \(H=G_{\underline M}\) and
\(K=\operatorname{an}_G(M)\). Let \(E_2\) be the set of edges
\(u\to v\) of \(G\) with \(u\in M\) and \(v\in K\), and let \(H_2\)
be \(H\) with the edges of \(E_2\) added back. It is a subgraph of
\(G\), hence acyclic, with the same vertices and private roots. In
\(H_2\) a vertex of \(K\) has all of its parents from \(G\): an edge
into it from a mediator is in \(E_2\), and every other edge into it
is in \(H\). A vertex outside \(K\) has in \(H_2\) exactly its
non-mediator parents from \(G\). Moreover
\(K=\operatorname{an}_H(M)\). Indeed, for \(v\in K\setminus M\)
follow a directed route in \(G\) from \(v\) to a mediator and stop at
the first mediator on it. Every edge before that endpoint leaves a
non-mediator, so the shortened route lies in \(H\). The reverse
inclusion holds because \(H\) is a subgraph of \(G\).

\emph{The auxiliary models.} For each mediator value tuple \(m\),
keep the natural equation \(f_v(x_{\operatorname{pa}(v)},u_v)\) at
every vertex \(v\in K\), every mediator included. At a vertex
\(v\notin K\), keep \(f_v\) but feed the coordinate \(m_u\) into the
input position of each mediator parent \(u\). Each new function is the
composite of \(f_v\) with a measurable coordinate substitution, so it
is jointly measurable in \(m\), the remaining parent values and the
disturbance. The disturbances and their joint law are unchanged, the
graph of every auxiliary model is \(H_2\), and evaluation in a
topological order of \(H_2\) gives each auxiliary model its unique
solution, jointly measurable in \(m\) and the disturbances by
induction along that order.

\emph{The tuples \(M\) and \(T\) are unchanged.} Let \(N_0\) be the
union of \(K\) with the set of vertices that are not descendants of
any mediator. It is ancestral in \(G\): an ancestor of a vertex of
\(K\) lies in \(K\), and an ancestor of a vertex that does not descend
from a mediator does not descend from one either. At a vertex of
\(N_0\) the auxiliary equation is the natural equation with natural
inputs. For \(v\in K\) this holds by construction. A vertex of
\(N_0\setminus K\) has no mediator parent, since such a parent would
make it a mediator descendant, so nothing in its equation was
replaced. Induction along a topological order of \(N_0\) shows that
every vertex of \(N_0\) takes its actual value, at every outcome and
for every \(m\). The mediators lie in \(K\), and the coordinates of
\(T\) are not mediator descendants. Hence the auxiliary tuples
\(M\) and \(T\) coincide pointwise with the actual ones for every
\(m\).

\emph{No new ancestral vertex.} Put
\(\mathcal A=\operatorname{an}_H(\{Y\}\cup M\cup T)\). Then
\(\operatorname{an}_{H_2}(\{Y\}\cup M\cup T)=\mathcal A\). The
inclusion of \(\mathcal A\) holds because \(H\) is a subgraph of
\(H_2\). Conversely, let a directed route in \(H_2\) lead from \(x\)
to \(\{Y\}\cup M\cup T\). If it uses no edge of \(E_2\), it lies in
\(H\). Otherwise its first edge from \(E_2\) leaves a mediator \(u\),
and the part of the route before \(u\) is a directed route in \(H\)
from \(x\) to \(u\in M\). In both cases \(x\in\mathcal A\). Each edge
of \(E_2\) has both ends in \(\mathcal A\), because its tail is a
mediator and its head lies in \(K=\operatorname{an}_H(M)\subseteq
\mathcal A\). Thus the restored ancestral graph has the same vertex
set \(\mathcal A\) and differs from the old one only by the edges
\(E_2\).

\emph{New moral connections stay on the mediator side.} Write
\(\mathcal M\) and \(\mathcal M_2\) for the moral graphs
(\cref{def:p16-ancestral-moral-graph}) of \(H\) and \(H_2\)
restricted to \(\mathcal A\). Let \(D_Y\) be the component of \(Y\)
in \(\mathcal M\setminus T\), and let \(D_M\) be the union of the
components of \(\mathcal M\setminus T\) meeting \(M\). Exposing
private roots does not alter separation among the original vertices,
so the third separation and \cref{thm:p16-moral-separation-bridge}
give \(D_Y\cap D_M=\varnothing\). Moralization joins two vertices
only through a directed edge or a common child, and the only vertices
whose parent sets grew are heads of edges in \(E_2\). Hence every
edge of \(\mathcal M_2\) that is absent from \(\mathcal M\) has both
endpoints in \(\{v\}\cup\operatorname{pa}_{H_2}(v)\) for the head
\(v\) of some edge of \(E_2\).

Let \(u\to v\) be an edge of \(E_2\). Its head \(v\) lies in \(D_M\).
It is a child of the mediator \(u\), and since
\(v\in\operatorname{an}_H(M)\) there is a directed route in \(H\)
from \(v\) to a mediator. Every vertex on that route is a descendant
of \(u\) in \(G\), so none of them lies in \(T\), and all of them
lie in \(\mathcal A\). The edges of the route are edges of
\(\mathcal M\setminus T\), so there \(v\) is joined to a mediator.
Now let \(w\) be any other parent of \(v\) in \(H_2\) that is not in
\(T\). If \(w\) is a mediator, then \(w\in M\subseteq D_M\).
Otherwise \(w\to v\) is an edge of \(H\) with \(w\in\mathcal A\),
so \(w\) and \(v\) are adjacent in \(\mathcal M\setminus T\) and
\(w\in D_M\). This case includes the private root of \(v\). Thus
restoring the edge puts its head, and every parent of that head
outside \(T\), on the mediator side.

Delete \(T\) from \(\mathcal M_2\). An added edge with an endpoint in
\(T\) disappears, and every other added edge joins two vertices of
\(D_M\). No added edge therefore has an endpoint in \(D_Y\). The edges
of \(\mathcal M\setminus T\) inside \(D_Y\) are still present, so
\(D_Y\) is again a connected component of \(\mathcal M_2\setminus T\),
and it does not meet \(M\). By \cref{thm:p16-moral-separation-bridge}
applied to \(H_2\), the vertex \(Y\) is d-separated from \(M\) given
\(T\) in \(H_2\). The component used by the response-component lemma
for \(Y\), the set \(B=M\) and the conditioning set \(T\) in \(H_2\)
is exactly \(D_Y\).

\emph{The auxiliary response is \(Y_m\).} Compare the auxiliary model
at \(m\) with the model in which the mediator equations are replaced
by the constants \(m\) (\cref{def:p16-measurable-scm}). Let \(x\) be
a non-root vertex of \(D_Y\). It is not a mediator, since
\(M\subseteq D_M\) and \(D_Y\cap D_M=\varnothing\). If \(x\in K\) had
a mediator parent in \(G\), that edge would belong to \(E_2\) and
\(x\) would lie in \(D_M\), which is impossible. Hence either \(x\)
has no mediator parent, or \(x\notin K\). In both cases the auxiliary
equation at \(x\) is \(f_x\) with \(m_u\) fed into the position of
each mediator parent \(u\). This is exactly its equation in the
intervened model, where \(u\) is the constant \(m_u\). Every
non-mediator parent of \(x\), its private root included, is joined to
\(x\) by an edge of \(H\) and lies in \(\mathcal A\), so it lies in
\(D_Y\) or in \(T\). In both models the coordinates of \(T\) take
their actual values. For the auxiliary model this was shown above. For
the intervened model, the set of vertices that are neither mediators
nor mediator descendants is ancestral, since an ancestor that is a
mediator would make its descendants mediator descendants; none of its
equations was replaced, and it contains every coordinate of \(T\).
The disturbances are the same in both models.
Induction along a topological order of \(D_Y\) then shows that every
vertex of \(D_Y\) takes the same value in the two models, at every
outcome and for every \(m\). In particular the auxiliary response at
\(m\) equals \(Y_m\) pointwise.

\emph{Conclusion.} The auxiliary models form a family on the single
DAG \(H_2\), with the same independent disturbances and a standard
Borel parameter \(m\) entering jointly measurable functions. The
tuples \(B=M\) and \(T\) are disjoint, exclude \(Y\), and are
pointwise unchanged as \(m\) varies, and \(Y\) is d-separated from
\(M\) given \(T\) in \(H_2\). The response-component lemma therefore
gives the tuple \(R_2\) of private roots in \(D_Y\) and one jointly
measurable \(F_2\) with \(Y_m=F_2(m,T,R_2)\) pointwise for every
\(m\), and \(R_2\perp M\mid T\). The auxiliary \(M\), \(T\) and roots
are the actual ones pointwise, so this independence is a statement
about the original joint law. In particular consistency supplies the
simultaneous measurable identity \(Y=F_2(M,T,R_2)\). This identity,
not conditioning on a null equality event between two continuously
distributed mediators, is what will identify the observed response
kernel.

\emph{A worked boundary.} Take mediators \(M=(M_1,M_2)\), a further
vertex \(Z\), empty covariate sets, and edges
\(A\to M_1\to Z\to M_2\to Y\), \(M_1\to Y\) and \(A\to Y\), so
that \(T=(A)\). In \(G_{\underline M}\) the only edges at \(Y\) come
from \(A\) and from the private root of \(Y\), and \(A\in T\), so the
third separation holds. Here \(E_2\) consists of the single edge
\(M_1\to Z\); the edges \(M_1\to Y\) and \(M_2\to Y\) are not
restored, since \(Y\notin K\). The auxiliary model evaluates \(Z\)
from the natural \(M_1\) and \(M_2\) from the natural \(Z\), while
the response is \(f_Y(A,m_1,m_2,U_Y)=Y_m\). The set \(D_Y\) consists
of \(Y\) and its root, and restoration adds the moral edges joining
\(M_1\) to \(Z\) and to the root of \(Z\), both inside \(D_M\). If
the edge \(Z\to Y\) is added, the response \(Y_m\) needs \(Z\)
evaluated at \(m_1\), whereas the restored model evaluates \(Z\) at
the natural \(M_1\). That graph also violates the third separation:
the path \(Y\leftarrow Z\to M_2\) lies in \(G_{\underline M}\), and
its middle vertex is a non-collider descending from \(M_1\), hence
outside \(T\).

\paragraph{The reference-mediator law.}
Finally retain natural \(A\) and replace its outgoing inputs by
\(b\). The mediator is \(M_b\), and \((C,W_3)\) is unchanged.
The fourth separation and the general global Markov theorem give
\(M_b\perp A\mid C,W_3\). Positive conditional exposure atoms
identify its regular conditional law there with the observed law
of \(M\) in exposure group \(b\), using consistency. The equality
is outside one \((C,W_3)\)-null set as measures: compare a countable
generating class and use measure uniqueness. Integrate over
\(\gamma_3\) to obtain the displayed \(\eta_b\) as the law of
\(M_b\) given \(C\), again with one common baseline-null set.
All three field independences have this same common-measure
interpretation by the conditional-independence calculus.
\end{proof}
\end{annexitem}

\begin{annexitem}{annex:p16.measurable-mediation-corollary}{thm:p16-measurable-mediation-corollary}{Independent reference draw, sequential replacement and overlap transfer}
\begin{proof}
\paragraph{Replace the natural coupling by an independent draw using its proved field.}
Construct a new mediator \(G^b\) with joint law
\(P(d\omega)\eta_b(C(\omega),dg)\), as in the preceding mediation
policy definition. This uses the original background, together with a
fresh draw conditional on baseline. The kernel joint-measure theorem
proves existence. Conditional on \(C\), the joint law of
\((M_b,R_0)\) is the product of \(\eta_b\) and the conditional
\(R_0\) law by the first field. The enlarged model has exactly that
same product law for \((G^b,R_0)\). The single jointly measurable
map \(F_0(a,\cdot,C,\cdot)\) therefore gives
\[
 \mathcal L(Y_{a,M_b})=\mathcal L(Y_{a,G^b}).               \tag{1}
\]
This is a statement about equal pushforwards of equal joint measures;
it involves no pointwise conditional-version substitution.

\paragraph{Use the exposure field on that independent draw.}
Because \(S\) contains \(C\), the conditional law of
\((R_1,A,G^b)\) given \(S\) factors as
\[
 \mathcal L(R_1\mid S)\otimes\mathcal L(A\mid S)
                         \otimes\eta_b(C,\cdot).
\]
To verify this, integrate a rectangle indicator under the enlarged
measure. The new draw contributes the factor \(\eta_b(C,D)\),
which is measurable from \(S\), and the first two factors follow
from \(R_1\perp A\mid S\). A countable generating-class argument
gives equality as kernels outside one \(S\)-null set. Consequently
\((R_1,G^b)\perp A\mid S\), and the joint response family gives
\(Y_{a,G^b}=F_1(a,G^b,S,R_1)\perp A\mid S\).
For a bounded \(\varphi\), positive \(P(A=a\mid S)\) now permits
\[
 E[\varphi(Y_{a,G^b})]
 =\int E[\varphi(Y_{a,G^b})\mid A=a,S=s]\,P_S(ds).         \tag{2}
\]
The baseline law of \(S=(C,W_1)\) is \(\nu\gamma_1\).
Its law given \(A=a\) has density
\(P(A=a\mid S)/P(A=a)\) with respect to it; this is a positive
Radon--Nikodym derivative almost everywhere, not an assumed coordinate
density. Thus the two laws are equivalent, and the conditional equality
in (2) holds on every needed measure class.

\paragraph{Use the mediator field before evaluating its random argument.}
Let \(\kappa_2(t,dr)\) be a conditional law of \(R_2\) given
\(T=t\). The third field gives one product kernel for
\((M,R_2)\) given \(T\). Since
\(Y=F_2(M,T,R_2)\) pointwise, the jointly measurable kernel
\[
 \widetilde Q(t,m,B)=\int\mathbf1_B(F_2(m,t,r))\kappa_2(t,dr)
\]
is a conditional law of the observed \(Y\) given \((T,M)\).
This follows by integrating the actual product law on rectangles;
kernel uniqueness extends the equality to all response events on
one \(P_{T,M}\)-null set. In exposure group \(a\), it therefore
agrees with the specified \(Q_a^{\mathrm{obs}}\) for the observed
law of \((C,W_1,W_2,M)\) in that group.

The new draw has law \(\eta_b(C,\cdot)\) independent of \(R_2\)
conditional on \(T\): the same rectangle calculation applies because
\(T\) includes \(C\). Thus a conditional law of \(Y_{G^b}\) given
\(T=t\) is \(\int\widetilde Q(t,m,\cdot)\eta_b(c,dm)\).
On the event \(A=a\), intervention on \(M\) leaves \(A\) unchanged,
since every mediator is its descendant. Composition at that attained
exposure value gives \(Y_{a,G^b}=Y_{G^b}\) pointwise in the enlarged
space. Partitioning the right side of (2) by \(W_2\) consequently
gives the target integral with \(\widetilde Q\), \(\eta_b\),
\(\gamma_2^a\), \(\gamma_1\), and \(\nu\).

It remains to justify replacement of \(\widetilde Q\) by the
\emph{observed} response kernel on this new integration measure.
Let \(\alpha=\nu\gamma_1\gamma_2^a\). The observed group law
of \((C,W_1,W_2,M)\) is equivalent to \(\alpha\rho_a\),
by the positive exposure derivative above. If a measurable set has
zero measure under this law, nonnegative kernel integration implies
that its mediator section has \(\rho_a\)-mass zero for
\(\alpha\)-almost every \((c,w_1,w_2)\). Condition (MS) gives
zero \(\eta_b\)-mass for all those sections. Therefore the
target measure \(\alpha\eta_b\) is absolutely continuous with
respect to the observed group law. The single exceptional set on
which the two conditional response measures differ has target mass
zero. Their integrals are identical, proving the stated observable
formula together with (1).

For means, first apply the bounded formula to \(|y|\wedge n\).
Monotone convergence identifies the target absolute integral with
the absolute expectation of the nested response. Under its stated
finiteness, positive and negative parts give the signed formula.
The observed integrability condition permits using the ordinary
conditional mean version, with a measurable default on the
observed-null infinite-integral set; absolute continuity also removes
that set from the target. Indicator tests identify the entire response
law. Applying this result to a finite collection of exposure pairs
and taking a finite union of the relevant null sets licenses all
the stated effect subtractions.
\end{proof}
\end{annexitem}

\begin{annexitem}{annex:p16.measurable-graph-examples}{thm:p16-measurable-graph-examples}{Singularity does not prevent kernel factorization}
\begin{proof}
\paragraph{The singular chain: kernels and the global property.}
Give each vertex of \(X\to Z\to Y\) the coordinate space
\(\mathbb R\). The chain model has kernels \(K_X\) uniform on
\((0,1)\), \(K_Z(x,\cdot)=\delta_x\) and \(K_Y(z,\cdot)=\delta_z\).
The point-mass kernel \(x\mapsto\delta_x(B)=\mathbf1_B(x)\) is
measurable for every Borel \(B\), so these are probability kernels
specified at every parent argument. The iterated construction of
\cref{def:p16-measurable-dag-kernels} gives the law of
\((U,U,U)\), and \cref{thm:p16-measurable-global-markov} makes it
globally Markov.

\paragraph{The singular chain is not product dominated.}
Its law gives
probability one to \(D=\{(x,z,y):x=z=y\}\). Suppose a product
\(\mu_1\otimes\mu_2\otimes\mu_3\) of sigma-finite measures
dominated this law. A sigma-finite measure on a real Borel space has
at most countably many positive-mass singletons: cover the space by
countably many finite-measure sets, and in each such set there are
only finitely many singletons with mass at least \(1/n\), for each
positive integer \(n\). Let \(J\) be the countable set of points with
positive \(\mu_2\)-mass. Every singleton has finite mass under a sigma-finite
measure, since it lies in one member of such a finite-measure cover.
For \(x\notin J\), the section of
\(D\cap\{x\notin J\}\) has product mass
\(\mu_2(\{x\})\mu_3(\{x\})=0\). Tonelli therefore gives this
whole diagonal subset product measure zero. Its actual probability
is one, since a uniform variable avoids the countable \(J\).
This contradicts absolute continuity. Thus a product-density proof
cannot be presumed for these kernels.

\paragraph{The singular chain: every fixed structural intervention.}
Write the chain as a measurable structural model in the sense of
\cref{def:p16-measurable-scm}: the disturbance of \(X\) is \(U\),
uniform on the standard Borel space \((0,1)\); the disturbances of
\(Z\) and \(Y\) take values in a one-point space \(\{\ast\}\); and
\[
 f_X(u)=u,\qquad f_Z(x,\ast)=x,\qquad f_Y(z,\ast)=z .
\]
Each function is a coordinate projection, hence jointly measurable and
total. Let \(I\subseteq\{X,Z,Y\}\) be any of the eight intervention
sets and \(t\in\mathbb R^I\). Solving in the order \(X,Z,Y\) gives
the only possible solution
\[
 X^t=\begin{cases}t_X,&X\in I,\\U,&X\notin I,\end{cases}\quad
 Z^t=\begin{cases}t_Z,&Z\in I,\\X^t,&Z\notin I,\end{cases}\quad
 Y^t=\begin{cases}t_Y,&Y\in I,\\Z^t,&Y\notin I.\end{cases}
\]
In each case the value is a coordinate of \((t,U)\), so the solution
is jointly measurable in \((t,U)\). Its law is the iterated kernel
construction in which the kernel of each \(v\in I\) is replaced by
\(\delta_{t_v}\) and the others are kept: at \(X\) the factor is
either \(\delta_{t_X}\) or the uniform law of \(U\), and at \(Z\) and
\(Y\) it is either \(\delta_{t_v}\) or the point mass at the parent
value, which is the value the equation copies. This is the conclusion
of \cref{thm:p16-measurable-interventions}, checked directly for every
\(I\) and \(t\). For instance, the intervention \(Z=z\) has solution
\((U,z,z)\) and law
\(\operatorname{Unif}(0,1)\otimes\delta_z\otimes\delta_z\) for every
real \(z\), although \(P(Z=z)=P(U=z)=0\) for every \(z\). So all fixed
structural interventions exist, including those whose prescribed value
is an observed null event.

\paragraph{The mediation example: hypotheses of the corollary.}
The structural model on the graph \(A\to M\to Y\), \(A\to Y\) has
mutually independent disturbances: the fair bit for \(A\), the uniform
\(U\) for \(M\), and a one-point disturbance for \(Y\). Its equations
are the identity on the bit for \(A\),
\(f_M(a,u)=u\,\mathbf1_{\{a=0\}}+\sqrt u\,\mathbf1_{\{a=1\}}\) and
\(f_Y(a,m)=a+m+am\), all jointly measurable. The single mediator
\(M\) is a child of \(A\) and a parent of \(Y\). All covariate sets
are empty, so \(C\) is trivial, \(\nu\) is a point mass, and
\(\gamma_1,\gamma_3,\gamma_2^a\) are point masses on the empty tuple.

The four separations of \cref{thm:p16-mediation-four-set-criterion}
are checked edge by edge. In \(G_{\underline A\underline M}\) every
arrow leaves \(A\) or \(M\), so the graph has no edges and \(Y\) is
separated from \(M\) given the empty set. In
\(G_{\underline A\overline M}\) the arrows \(A\to M\) and \(A\to Y\)
are deleted, leaving only \(M\to Y\); \(A\) is isolated, so \(Y\) is
separated from \(A\). In \(G_{\underline M}\) only \(M\to Y\) is
deleted, leaving \(A\to M\) and \(A\to Y\); the single path
\(M\leftarrow A\to Y\) has the non-collider \(A\) in the conditioning
set, so \(Y\) is separated from \(M\) given \(A\). In
\(G_{\underline A}\) only \(M\to Y\) remains; \(A\) is isolated, so
\(M\) is separated from \(A\). The nondescendant requirements on the
empty covariate sets hold vacuously.

The exposure atoms are \(P(A=0)=P(A=1)=1/2>0\). By independence of
\(A\) and \(U\), \(\mathcal L(M\mid A=0)\) is the law of \(U\) and
\(\mathcal L(M\mid A=1)\) is the law of \(\sqrt U\). For \(0<m<1\),
\(P(U\leq m)=m\) and \(P(\sqrt U\leq m)=P(U\leq m^2)=m^2\), so these
laws have Lebesgue densities \(1\) and \(2m\) on \((0,1)\). Thus
\(\rho_a\) is the law with density \(1\) or \(2m\) according as
\(a=0\) or \(a=1\), and with empty \(W_3\) the reference kernel
\(\eta_b\) is the law with density \(1\) or \(2m\) according as
\(b=0\) or \(b=1\). Both densities are positive on \((0,1)\), so each
law gives zero mass to exactly the Lebesgue-null subsets of
\((0,1)\); hence \(\eta_b\ll\rho_a\) for all four pairs \((a,b)\),
which is (MS). Since \(Y=A+M+AM\) pointwise,
\(Q^{\mathrm{obs}}_a(m,\cdot)=\delta_{a+m+am}\) is a version of
\(\mathcal L(Y\mid A=a,M=m)\): for Borel sets \(B_1,B_2,B_3\),
\(P(A\in B_1,M\in B_2,Y\in B_3)=
E[\mathbf1_{B_1}(A)\mathbf1_{B_2}(M)\delta_{A+M+AM}(B_3)]\).
The response lies on the curve \(y=a+m+am\) and has no ordinary
conditional density on a response interval; the corollary uses only
this conditional probability kernel.

\paragraph{The mediation example: integrability and the four values.}
For \(a\in\{0,1\}\) and \(0<m<1\), \(0\leq a+m+am<3\). Hence
\(0\leq Y<3\), so \(E|Y|\leq3<\infty\), and the absolute response
integral of the corollary is
\[
 \int\!\int|y|\,Q^{\mathrm{obs}}_a(m,dy)\,\eta_b(dm)
 =\int_0^1|a+m+am|\,\eta_b(dm)\leq3<\infty
\]
for each of the four pairs. Both mean premises hold, so the mean form
of \cref{thm:p16-measurable-mediation-corollary} gives
\(E[Y_{a,M_b}]=\int_0^1(a+m+am)\,\eta_b(dm)\). With \(a_0=0,a_1=1\)
and \(\theta_{ij}=E[Y_{a_i,M_{a_j}}]\) from
\cref{def:p16-mediation-effects},
\[
 \theta_{00}=\int_0^1m\,dm=\tfrac12,\qquad
 \theta_{01}=\int_0^1m\cdot2m\,dm=\tfrac23,
\]
\[
 \theta_{10}=\int_0^1(1+2m)\,dm=2,\qquad
 \theta_{11}=\int_0^1(1+2m)\,2m\,dm=1+\tfrac43=\tfrac73.
\]
As a structural check, \(Y_{a,M_b}=a+(1+a)M_b\) with \(M_0=U\),
\(M_1=\sqrt U\), \(E[U]=1/2\) and \(E[\sqrt U]=2/3\) give the same four
numbers. Subtraction gives
\(\operatorname{PNDE}=\theta_{10}-\theta_{00}=3/2\),
\(\operatorname{TNIE}=\theta_{11}-\theta_{10}=1/3\), and
\(\operatorname{TE}=\theta_{11}-\theta_{00}=11/6\).

\paragraph{The null exposure point: two structural models.}
Let \(A\), \(M\) and \(E\) be independent, with \(A\) and \(M\)
uniform on \((0,1)\) and \(E\) a Bernoulli\((1/4)\) bit, and fix
\(a_*=1/2\), \(b_*=1/4\). Both models use the graph
\(A\to M\to Y\), \(A\to Y\), the disturbances \(A\), \(M\), \(E\)
of the three vertices, and \(f_A(u)=u\), \(f_M(a,u)=u\) (the
permitted input \(a\) is unused). The response equations are
\[
 f^{(1)}_Y(a,m,e)=e\oplus\mathbf1_{\{a=1/2\}},\qquad
 f^{(2)}_Y(a,m,e)=e .
\]
Both are jointly measurable, since \(\{a=1/2\}\) is a Borel set. The
graph is the one just checked, so the four separations hold with
empty covariate sets.

\paragraph{The two observed laws coincide.}
The observed variables are \((A,M,Y)\). In the second model \(Y=E\)
identically. In the first model \(Y\neq E\) only on \(\{A=1/2\}\),
and \(P(A=1/2)=0\) because \(A\) is uniform. Hence \(Y=E\) almost
surely in the first model as well. In both models, therefore,
\((A,M,Y)=(A,M,E)\) almost surely, and the observed law of
\((A,M,Y)\) is the product of two uniform laws on \((0,1)\) and the
Bernoulli\((1/4)\) law. Its density with respect to Lebesgue measure
in \((A,M)\) and counting measure on \(\{0,1\}\) for \(Y\) is \(3/4\)
at \(y=0\) and \(1/4\) at \(y=1\), at every point of the open unit
square; it is strictly positive. In both models \(M\) is independent
of \(A\), so the conditional mediator density is \(1\) at every
exposure value. Every functional of the observed law, including
every version of every observed conditional kernel, is thus common to
the two models.

\paragraph{The two natural direct effects differ.}
Intervening on \(A\) leaves \(M\) unchanged, so \(M_{b_*}=M\). In
the first model \(Y_{a,m}=E\oplus\mathbf1_{\{a=1/2\}}\), so
\(Y_{a_*,M_{b_*}}=1-E\) and \(Y_{b_*,M_{b_*}}=E\); their means are
\(3/4\) and \(1/4\), and the natural direct effect from \(b_*\) to
\(a_*\) is \(3/4-1/4=1/2\). In the second model both responses equal
\(E\), and the effect is zero. Both responses are bounded, so the
means exist. Since the two models have one observed law but
different effects, no functional of the observed law equals the
natural direct effect at the fixed exposure \(a_*\) in both. The
corollary's atom premise fails here, since \(P(A=a_*)=0\). The
structural intervention is defined in either model, but a
conditional response value at the null exposure is not identified
by the observed law. Thus the conditional-atom domain in the
observational corollary is a stated sufficient domain; replacing
it by an unqualified point-density condition would be false.
\end{proof}
\end{annexitem}
